\section{The actual canonical ratio section is a frame off both supports}

\begin{lemma}[An invertible pulled functional has an invertible transpose]
\label{lem:original-pulled-transpose-frame}
\lean{EnvelopeProperty.PLOriginalModelCompatibility.RatioSectionTransport.pulledDualTranspose_isIso_of_functional_isIso}
\uses{lem:original-pl-ratio-pulled-functional-transport}
\leanok
Let $M$ be an invertible module on $N$, let
$a:M\to\mathcal O_N$ be an original functional, and let $f:T\to N$.
If the pulled functional $f^*M\to\mathcal O_T$ is an isomorphism,
then the pulled original evaluation transpose
$\mathcal O_T\to f^*(M^\vee)$ is an isomorphism.
\end{lemma}
\begin{proof}
Tensor the pulled transpose by $f^*M$ and compose with the actual
pulled evaluation isomorphism. The pulled transpose pairing identity
expresses this composite as the canonical right-unit comparison,
followed by the pulled functional and the canonical structure-sheaf
unit comparison. Every factor is an isomorphism. Cancel the actual
evaluation isomorphism, so tensoring the transpose by $f^*M$ gives
an isomorphism. Tensoring by this invertible module is fully faithful,
and therefore reflects isomorphisms. Hence the transpose itself is
an isomorphism.
\end{proof}

\begin{theorem}[The canonical ratio is its own actual frame]
\label{thm:original-canonical-ratio-section-frame}
\lean{EnvelopeProperty.PLOriginalModelCompatibility.RatioSectionTransport.canonicalPulledRatioSection_isIso, EnvelopeProperty.PLOriginalModelCompatibility.RatioSectionTransport.localSection_isIso}
\uses{lem:original-pulled-transpose-frame,thm:original-pl-canonical-ratio-section-transport}
\leanok
For original ideals $I,J$ with $J$ invertible, suppose both actual
pulled inclusions $f^*I\to\mathcal O_T$ and
$f^*J\to\mathcal O_T$ are isomorphisms. Then the original canonical
ratio section $\mathcal O_T\to f^*(I\otimes J^\vee)$ is an
isomorphism. In particular the actual original local ratio section
is a frame on every open disjoint from both ideal supports.
\end{theorem}
\begin{proof}
The numerator part is the inverse of the actual pulled numerator
inclusion and hence is an isomorphism. The preceding lemma shows
that the denominator part, the pullback of the original denominator
functional's evaluation transpose, is an isomorphism as well.
The tensor of these section maps, preceded by the canonical
structure-sheaf multiplication inverse, is therefore an isomorphism.
The actual coherent-to-explicit tensor comparison and canonical
pullback tensor comparison are isomorphisms, so composing with
them gives the stated canonical ratio section isomorphism. Off
both supports the two pulled inclusions are the actual canonical
open inclusion frames. The proved equality with the
original local ratio section gives the last assertion.
\end{proof}
